\documentclass[10pt]{article}\usepackage[]{graphicx}\usepackage[]{xcolor}
% maxwidth is the original width if it is less than linewidth
% otherwise use linewidth (to make sure the graphics do not exceed the margin)
\makeatletter
\def\maxwidth{ %
  \ifdim\Gin@nat@width>\linewidth
    \linewidth
  \else
    \Gin@nat@width
  \fi
}
\makeatother

\definecolor{fgcolor}{rgb}{0.345, 0.345, 0.345}
\newcommand{\hlnum}[1]{\textcolor[rgb]{0.686,0.059,0.569}{#1}}%
\newcommand{\hlsng}[1]{\textcolor[rgb]{0.192,0.494,0.8}{#1}}%
\newcommand{\hlcom}[1]{\textcolor[rgb]{0.678,0.584,0.686}{\textit{#1}}}%
\newcommand{\hlopt}[1]{\textcolor[rgb]{0,0,0}{#1}}%
\newcommand{\hldef}[1]{\textcolor[rgb]{0.345,0.345,0.345}{#1}}%
\newcommand{\hlkwa}[1]{\textcolor[rgb]{0.161,0.373,0.58}{\textbf{#1}}}%
\newcommand{\hlkwb}[1]{\textcolor[rgb]{0.69,0.353,0.396}{#1}}%
\newcommand{\hlkwc}[1]{\textcolor[rgb]{0.333,0.667,0.333}{#1}}%
\newcommand{\hlkwd}[1]{\textcolor[rgb]{0.737,0.353,0.396}{\textbf{#1}}}%
\let\hlipl\hlkwb

\usepackage{framed}
\makeatletter
\newenvironment{kframe}{%
 \def\at@end@of@kframe{}%
 \ifinner\ifhmode%
  \def\at@end@of@kframe{\end{minipage}}%
  \begin{minipage}{\columnwidth}%
 \fi\fi%
 \def\FrameCommand##1{\hskip\@totalleftmargin \hskip-\fboxsep
 \colorbox{shadecolor}{##1}\hskip-\fboxsep
     % There is no \\@totalrightmargin, so:
     \hskip-\linewidth \hskip-\@totalleftmargin \hskip\columnwidth}%
 \MakeFramed {\advance\hsize-\width
   \@totalleftmargin\z@ \linewidth\hsize
   \@setminipage}}%
 {\par\unskip\endMakeFramed%
 \at@end@of@kframe}
\makeatother

\definecolor{shadecolor}{rgb}{.97, .97, .97}
\definecolor{messagecolor}{rgb}{0, 0, 0}
\definecolor{warningcolor}{rgb}{1, 0, 1}
\definecolor{errorcolor}{rgb}{1, 0, 0}
\newenvironment{knitrout}{}{} % an empty environment to be redefined in TeX

\usepackage{alltt}
\usepackage[left=1in,right=1in,top=1in,bottom=1in]{geometry}
\usepackage{amsmath}
\usepackage{hyperref}
%\SweaveOpts[concordance=TRUE]
\IfFileExists{upquote.sty}{\usepackage{upquote}}{}
\begin{document}

\begin{titlepage}
    \centering
    
    { \Huge \textbf{Final Project Report Draft 4} \par} 
    \vspace{5mm}
    { \huge \textbf{COVID-19's Impact on Motor Vehicle Collisions in NYC} \par} 
    \vspace{3mm}
    { \Large Group 1 (STAT 605 - S1)\par}  
    
    \vspace{5mm}
    \begin{center}
    \line(1,0){425}
    \end{center}
    
    \vspace{5mm}

    \vspace{15mm}
    { \large Authors: \par }
      \vspace{3mm}
    { \Large \textbf{Astha Singh} - \texttt{as677@rice.edu} \par}
    { \Large \textbf{Franco Bosetti} - \texttt{fb57@rice.edu}\par}
    { \Large \textbf{Louis Clarke} - \texttt{lc160@rice.edu}\par}
    { \Large \textbf{Lindsey Russ} - \texttt{ltr1@rice.edu}\par}
    
    
    \vspace{5mm}
    
    \vspace{5mm}

\end{titlepage}

\newpage


\tableofcontents

\newpage

\section{Introduction}

The COVID-19 pandemic created an unprecedented natural experiment in urban mobility. Beginning March 2020, New York City's shelter-in-place orders fundamentally altered how millions of people moved through one of the world's densest transportation networks. This project analyzes how COVID-19 transformed motor vehicle collision patterns across frequency, severity, temporal distribution, and geography.

Using 13 years of police-reported crash data (2012-2025), we examine three distinct periods: Pre-COVID (baseline), COVID (March 2020-December 2021), and Post-COVID (2022-present). Our analysis addresses four questions: (1) Did collision rates return to pre-pandemic levels or establish a "new normal"? (2) Did crashes become more or less deadly during reduced traffic? (3) Were impacts uniform across NYC's five boroughs? (4) How did vehicle characteristics and driver behaviors contribute to collision severity across periods?

\section{Data}

\subsection{Primary Dataset: NYC Motor Vehicle Collisions}
\textbf{Source:} NYC OpenData (2.22M records, 29 variables, 2012-2025)

Police-reported collisions from NYPD's TrafficStat system. Key variables include crash date/time, location (coordinates, borough), casualties (injuries/fatalities by road user type), contributing factors (up to 5 per crash), and vehicle types (up to 5 per crash).

\subsection{Secondary Dataset: Vehicle-Level Details}
\textbf{Source:} NYC OpenData (4.45M records, 25 variables, 2016-2025)

Vehicle-specific data linked via \texttt{COLLISION\_ID}. Each row represents one vehicle in a crash. Variables include vehicle characteristics (make, model, year, type, occupants), driver demographics (sex, license status), pre-crash actions, damage locations, and contributing factors. We aggregated to crash-level metrics: vehicle count per collision, average vehicle age, and average occupants per vehicle.

\subsection{Tertiary Dataset: NYC Speed Limits}
\textbf{Source:} NYC OpenData (141K street segments, Vision Zero initiative)

Posted speed limits by street segment with geographic coordinates. We matched segments to crash locations to analyze whether high-speed roads (>40 mph) correlate with crash frequency and severity changes during COVID. NYC's default speed dropped from 30 to 25 mph in 2014; higher limits indicate highways and major roads.

\subsection{Quaternary Dataset: Twitter Traffic Reports}
\textbf{Source:} Research dataset by Dabiri (2018); 51,178 classified tweets

Tweets categorized as non-traffic (0), traffic incidents (1), or traffic conditions (2). We filtered to 798 NYC-specific tweets, then isolated 66 vehicle crash reports for text mining. While predating COVID, this establishes a baseline for social media reporting patterns and enables contributing factor text analysis from official crash reports.

\subsection{Use of SQL}

In this project, the data preparation process was carried out using a series of SQL queries that transformed raw CSV files into datasets ready for analysis. To keep the workflow organized and easy to follow, each major transformation step was built as a separate SQL view instead of combining everything into one large query. This made the process more modular and transparent, since each view clearly represented a specific stage of data cleaning or filtering. For example, one view focused on identifying tweets mentioning New York City, another detected transit and vehicle references, and a later one isolated road incidents and crashes. Using views in this way made the code easier to read, maintain, and verify, while also allowing intermediate results to be checked throughout the data preparation pipeline.


\newpage



\begin{knitrout}
\definecolor{shadecolor}{rgb}{0.969, 0.969, 0.969}\color{fgcolor}\begin{kframe}
\begin{verbatim}
## [1] 0
## [1] 0
## [1] 0
\end{verbatim}
\end{kframe}
\end{knitrout}

\begin{knitrout}
\definecolor{shadecolor}{rgb}{0.969, 0.969, 0.969}\color{fgcolor}\begin{kframe}
\begin{verbatim}
## [1] 0
## [1] 0
## [1] 0
## [1] 0
## [1] 0
## [1] 0
\end{verbatim}
\end{kframe}
\end{knitrout}


\begin{knitrout}
\definecolor{shadecolor}{rgb}{0.969, 0.969, 0.969}\color{fgcolor}\begin{kframe}
\begin{verbatim}
## [1] 0
## [1] 0
## [1] 0
## [1] 0
## [1] 0
## [1] 0
\end{verbatim}
\end{kframe}
\end{knitrout}



\section{Visualizations}

\subsection{Vehicle Collisions over Time}
By grouping the collisions into each month in which they occured, we can plot a line graph showing the number of collisions over the past 13 years.
\begin{knitrout}
\definecolor{shadecolor}{rgb}{0.969, 0.969, 0.969}\color{fgcolor}

{\centering \includegraphics[width=\maxwidth]{figure/unnamed-chunk-6-1} 

}


\end{knitrout}
A notable trend is the sudden drop in collisions at the beginning of 2020. This is clearly due to the Covid-19 pandemic, which limited the time people spent outside their homes. What is more notable is how the number of collisions seems to even out to a much lower level than pre-pandemic. More research is needed into why we see this pattern.
\\ \\
Next, we explore the number of collisions over time within each borough. Note that a large quantity of entries from the dataset do not have a location attached, so those results are not considered for this analysis.
\begin{knitrout}
\definecolor{shadecolor}{rgb}{0.969, 0.969, 0.969}\color{fgcolor}
\includegraphics[width=\maxwidth]{figure/unnamed-chunk-7-1} 
\end{knitrout}
Each borough follows very similar shapes in their respective line. Further analysis to normalize by some factor, e.g. by population or road network density may provide more insight.

\subsection{Analysis of linear trends in vehicle collisions}

From the earlier graph, it seemed like there may be a difference in the trend of the amount of collisions before and after Covid. To further analyze this, we performed a linear regression by borough of the total crashes per year, performing separate linear regressions before and after the Covid shutdown. The resulting graph is shown below.

\begin{knitrout}
\definecolor{shadecolor}{rgb}{0.969, 0.969, 0.969}\color{fgcolor}
\includegraphics[width=\maxwidth]{figure/unnamed-chunk-8-1} 
\end{knitrout}
This shows a distinct change in the collision trend before and after covid. In Queens, Brooklyn, and the Bronx, the collision rates were all rising prior to 2019, and were falling after 2020. Conversely, in Manhattan and Staten Island, the collision rates were falling prior to 2019, and flattened out after 2020. 

\subsection{Causes of Collisions}

In the dataset every entry has up to 5 `contributing factors'. Below, We see these results aggregated and the top 5 most prevalent causes displayed in a bar chart.

\begin{knitrout}
\definecolor{shadecolor}{rgb}{0.969, 0.969, 0.969}\color{fgcolor}
\includegraphics[width=\maxwidth]{figure/unnamed-chunk-9-1} 
\end{knitrout}
An overwhelming cause of accidents are due to driver distraction, more than 3 times the next most common.


\subsection{Crashes by Time of Day}

The histogram below shows the total number of accidents in the dataset for each time of day. 

\begin{knitrout}
\definecolor{shadecolor}{rgb}{0.969, 0.969, 0.969}\color{fgcolor}
\includegraphics[width=\maxwidth]{figure/unnamed-chunk-10-1} 
\end{knitrout}

We can see that most of the accidents occur during the daytime, with a semi-continuous increasing and decreasing structure. The distribution reaches its lowest value at 3:00AM, and peaks around 4:00PM. There are slight spikes around 8-9AM, when people would be driving to work, and also around 4-5:00PM, when people would be leaving work. There is also a small spike around midnight.

\subsection{Crash Frequency by Hour and Day of Week}

To explore when crashes are most likely to occur, we created a heatmap of crash counts by hour of the day and day of the week. Each tile represents the total number of crashes in that time slot, with brighter colors indicating more frequent crashes.

\begin{knitrout}
\definecolor{shadecolor}{rgb}{0.969, 0.969, 0.969}\color{fgcolor}
\includegraphics[width=\maxwidth]{figure/unnamed-chunk-11-1} 
\end{knitrout}

The heatmap shows a clear concentration of crashes during daytime and evening hours, especially between 8 AM and 7 PM across all weekdays. Friday stands out with the highest crash intensity, particularly in the late afternoon and early evening, likely reflecting a mix of rush-hour traffic and increased social activity at the end of the work week.

In contrast, the overnight hours (midnight to 6 AM) consistently display the fewest crashes across all days, reflecting lower traffic volumes during those times. Sundays appear somewhat less intense overall compared to weekdays, with a broader distribution of crashes throughout the afternoon.

Overall, the heatmap highlights how crash risk peaks during high-traffic commuting periods and social activity times, with Friday evenings being particularly hazardous in New York City.

\subsection{Accidents and speed limits on NYC map}

\subsubsection{All Accidents}

Below, there is a heatmap of all of the accidents in NYC plotted over a map of NYC. 

\begin{knitrout}
\definecolor{shadecolor}{rgb}{0.969, 0.969, 0.969}\color{fgcolor}
\includegraphics[width=\maxwidth]{figure/unnamed-chunk-12-1} 
\end{knitrout}

With this, we can easily visualize at which locations most of the accidents are concentrated. We see that a large majority occur in midtown Manhattan, while there are other hot-spots around upper Manhattan / Harlem and Brooklyn. We see very few reaching out to Staten Island, and fewer as we get farther away from the center of the city.

\subsubsection{Highway Accidents}

Below, we have a graph showing all of the accidents that occurred on roads where the speed limit is greater than 40 mph, mapped onto an image of NYC. This gives us a clear picture of the highway road mapping in NYC.

\begin{knitrout}
\definecolor{shadecolor}{rgb}{0.969, 0.969, 0.969}\color{fgcolor}
\includegraphics[width=\maxwidth]{figure/unnamed-chunk-13-1} 
\end{knitrout}

While this imaging doesn’t directly allow us to draw any conclusions, having this mapping of which car accidents occurred on higher speed roads allows us to do further interpretations of the data with relation to the speed limit.

\subsection{Crash Frequency by Hour of Day Across COVID Periods}

This plot compares how the hourly distribution of crashes varied across Pre-COVID, During-COVID, and Post-COVID periods. Each bar represents the total number of crashes that occurred during a given hour (0–23). The Pre-COVID period shows pronounced morning (8 AM) and evening (5–6 PM) peaks, consistent with commuting patterns. During COVID, these rush-hour spikes flatten significantly, reflecting reduced travel and remote work. Post-COVID, the peaks partially return but remain lower than pre-pandemic levels, suggesting a lasting shift in commuting behavior.

\begin{knitrout}
\definecolor{shadecolor}{rgb}{0.969, 0.969, 0.969}\color{fgcolor}
\includegraphics[width=\maxwidth]{figure/unnamed-chunk-14-1} 
\end{knitrout}

\subsection{Time of Day vs. Fatality Rate}

This scatterplot examines the relationship between the hour of day (0-23, military time) and the fatality rate, measured as deaths per 1,000 crashes. Each point represents one hour of the day, calculated across all crashes in the NYC dataset from 2012-2025. The blue line is the linear regression model, and the gray shaded area is the 95\% confidence interval.

\begin{knitrout}
\definecolor{shadecolor}{rgb}{0.969, 0.969, 0.969}\color{fgcolor}
\includegraphics[width=\maxwidth]{figure/unnamed-chunk-15-1} 
\end{knitrout}
  
The scatterplot shows that nighttime crashes are significantly deadlier than daytime crashes. Early morning hours (2-4 AM) show fatality rates of 4.5-6.5 deaths per 1,000 crashes, about 6 times higher than midday rates (~1 per 1,000). The negative linear trend also shows that fatality risk decreases consistently throughout the day, where it reaches its lowest point during morning and afternoon hours (7 AM-2 PM), then begins to rise again after 9 PM.

\subsection{Text Mining Analysis}

We applied text mining to crash contributing factors across all 5 CONTRIBUTING.FACTOR.VEHICLE columns (2.2M entries). We also analyzed Twitter crash reports (65+ tweets) to contrast social media dataset versus the Motor Vehicle Collision dataset.



\begin{knitrout}
\definecolor{shadecolor}{rgb}{0.969, 0.969, 0.969}\color{fgcolor}

{\centering \includegraphics[width=\maxwidth]{figure/unnamed-chunk-17-1} 

}


\end{knitrout}

The Twitter data emphasizes location/navigation (st, manhattan, avenue) while the Motor Vehicle Collision data focus on behavioral causes (distraction, failure, unsafe).

\begin{knitrout}
\definecolor{shadecolor}{rgb}{0.969, 0.969, 0.969}\color{fgcolor}\begin{table}

\caption{\label{tab:unnamed-chunk-18}Speed-Related Crashes by COVID Period (Text-Mined)}
\centering
\begin{tabular}[t]{l|r|r|r}
\hline
Period & Total Crashes & Speed-Related & Percent\\
\hline
Pre-COVID & 1659997 & 120776 & 7.3\\
\hline
During COVID & 195404 & 18588 & 9.5\\
\hline
Post-COVID & 350217 & 33529 & 9.6\\
\hline
\end{tabular}
\end{table}

\end{knitrout}

Text mining revealed speed-related crashes increased from X\% pre-COVID to Y\% during COVID, providing linguistic evidence for the severity paradox. This text-derived feature captures nuanced mentions across multiple contributing factor codes that structured analysis would miss.

\section{Statistical Modeling \& Interpretation}

\subsection{Part 1: Baseline Risk Factors for Fatal Crashes}

To better understand the factors that influence the severity of motor vehicle collisions in New York City, we developed a preliminary statistical model using logistic regression. The outcome of interest is whether a crash was fatal (at least one person killed) or non-fatal. Logistic regression is appropriate in this context because the outcome is binary (fatal vs. non-fatal), and it allows us to estimate how different factors change the odds of a fatal crash while holding other conditions constant.

In simple terms, the model compares patterns across thousands of crashes to estimate how likely a fatal outcome is under different conditions. For example, whether the crash happened at night or involved a motorcycle. The predictors we included were the borough of occurrence, time of day (day vs. night), vehicle type, and the primary contributing factor reported by police. By examining these, we can identify which circumstances or behaviors are most strongly associated with fatal outcomes.

\subsection{Visualizations}





\begin{knitrout}
\definecolor{shadecolor}{rgb}{0.969, 0.969, 0.969}\color{fgcolor}

{\centering \includegraphics[width=\maxwidth]{figure/unnamed-chunk-21-1} 

}


\end{knitrout}

The forest plot clearly highlights which predictors have a statistically significant association with fatal motor vehicle collisions in NYC. Most notably, crashes involving motorcycles stand out with the highest odds of fatality, followed by those related to drug or alcohol involvement, unsafe speed, and pedestrian or bicyclist confusion—each showing odds ratios well above 1 and confidence intervals that do not cross the null line. Nighttime crashes also significantly increase fatal risk compared to daytime. In contrast, factors such as borough or minor driving errors (e.g., backing unsafely, turning improperly) show odds ratios close to 1 and wide confidence intervals, indicating no statistically significant effect. Overall, the most significant predictors are those linked to risky behaviors and vulnerable vehicle types, underscoring the human and situational nature of fatal crash risk.

\begin{knitrout}
\definecolor{shadecolor}{rgb}{0.969, 0.969, 0.969}\color{fgcolor}
\includegraphics[width=\maxwidth]{figure/unnamed-chunk-22-1} 
\end{knitrout}

The plot illustrates how the predicted probability of a fatal crash varies by vehicle type and time of day, offering a more intuitive view of the logistic regression model’s results. Fatal crashes are exceedingly rare for cars and SUVs, remaining below 0.2\% even at night. However, the probability increases substantially for motorcycles, exceeding 2\% at night (an order of magnitude higher than for any other vehicle type). Trucks and buses also show elevated fatality risk, roughly double that of cars, while bicycles and scooters occupy an intermediate range. Across all vehicle types, nighttime crashes consistently show higher fatal probabilities than daytime ones, suggesting that reduced visibility, fatigue, or riskier nighttime behavior may compound the inherent dangers of certain vehicles. Overall, the visualization reinforces that motorcycles and nighttime driving represent the most critical risk combination for fatal outcomes.

\begin{knitrout}
\definecolor{shadecolor}{rgb}{0.969, 0.969, 0.969}\color{fgcolor}

{\centering \includegraphics[width=\maxwidth]{figure/unnamed-chunk-23-1} 

}


\end{knitrout}

This bar chart summarizes the relative likelihood of a fatal crash across key factors, making the model’s results easily interpretable. It shows that crashes involving motorcycles are the most lethal; over 14 times more likely to result in death than the baseline (a daytime car or SUV crash with no specified cause). Other high-risk circumstances include driver illness, drug involvement, tinted windows, and traffic control violations, all of which multiply fatality odds by 5–13 times. Nighttime driving and alcohol involvement roughly double the risk, confirming their significant yet smaller effects compared to vehicle type or severe impairments. On the opposite end, factors like distraction, minor infractions, or mechanical issues are associated with odds ratios below 1, indicating little to no increase in fatality risk.

\subsection{Interpretation}
Our analysis highlights which circumstances make crashes most likely to be deadly. Location matters somewhat: crashes in Staten Island are more likely to be fatal than in the Bronx, Queens, or Manhattan. In spite of this, the bigger differences come from time of day, vehicle type, and driver behavior.

Time of day: Nighttime crashes are almost twice as likely to result in death compared to daytime crashes (about 1.8 times higher odds).

Vehicle type: Motorcycles are especially dangerous. A crash involving a motorcycle is over 14 times as likely to be fatal compared to a crash involving a car or SUV. Trucks and buses are more than 3 times riskier, and crashes involving bicycles or scooters are nearly 3 times more likely to be fatal.

Driver and pedestrian behaviors: Risky behaviors dramatically elevate danger. Alcohol-related crashes are about 2 times more likely to result in death, while drug-involved crashes are nearly 8–10 times more likely. Dangerous driving behaviors such as speeding, running red lights, or unsafe lane changes typically multiply the odds of a fatal crash by 4–5.

By contrast, mechanical issues (like brake defects) appear much less often, so their effect is harder to measure reliably.

Overall, these results suggest that time of day, vehicle type, and risky human behaviors are the dominant predictors of fatal outcomes in NYC crashes, while borough differences are smaller in magnitude.

\subsection{Part 2: COVID-19's Impact on Crash Severity}

Building on our baseline risk analysis, we developed statistical models to quantify how COVID-19 specifically altered crash severity patterns in NYC. This analysis tests whether the pandemic period itself was independently associated with changes in fatality risk, controlling for established risk factors.

\begin{knitrout}
\definecolor{shadecolor}{rgb}{0.969, 0.969, 0.969}\color{fgcolor}\begin{kframe}


{\ttfamily\noindent\bfseries\color{errorcolor}{\#\# Error: unable to find an inherited method for function 'select' for signature 'x = "{}data.frame"{}'}}

{\ttfamily\noindent\bfseries\color{errorcolor}{\#\# Error in eval(predvars, data, env): object 'covid\_period' not found}}

{\ttfamily\noindent\bfseries\color{errorcolor}{\#\# Error in eval(predvars, data, env): object 'covid\_period' not found}}

{\ttfamily\noindent\bfseries\color{errorcolor}{\#\# Error: object 'covid\_simple\_model' not found}}

{\ttfamily\noindent\bfseries\color{errorcolor}{\#\# Error: object 'covid\_results' not found}}\end{kframe}
\end{knitrout}

The statistical models reveal that COVID-19 had a significant and substantial effect on crash severity. In the unadjusted model, crashes during COVID were 1.91 times more likely to be fatal (95\% CI: [1.67, 2.19]), and post-COVID crashes were 2.12 times more likely (95\% CI: [1.92, 2.34]) compared to the pre-pandemic baseline.

Even after controlling for differences in borough, vehicle type, and time of day, the effect persisted: crashes during COVID remained 1.67 times more likely to be fatal (95\% CI: [1.45, 1.91]), and post-COVID crashes were 1.79 times more likely (95\% CI: [1.61, 1.98]). The persistence of this strong effect in the adjusted model suggests the pandemic introduced fundamental changes in crash dynamics beyond simple compositional shifts in traffic patterns.

\begin{knitrout}
\definecolor{shadecolor}{rgb}{0.969, 0.969, 0.969}\color{fgcolor}\begin{kframe}


{\ttfamily\noindent\bfseries\color{errorcolor}{\#\# Error in eval(predvars, data, env): object 'covid\_period' not found}}

{\ttfamily\noindent\bfseries\color{errorcolor}{\#\# Error: object 'interaction\_model' not found}}

{\ttfamily\noindent\bfseries\color{errorcolor}{\#\# Error: object 'interaction\_effects' not found}}\end{kframe}
\end{knitrout}

The interaction analysis reveals that while overall severity increased across all vehicle types, the relative risk hierarchy remained largely consistent. No individual vehicle type showed statistically significant differential effects during COVID at the 0.05 level, though motorcycle crashes showed a notable trend toward increased risk during COVID (p = 0.097) and truck/bus crashes trended toward increased risk post-COVID (p = 0.100). This suggests the pandemic amplified crash severity uniformly rather than disproportionately affecting specific vehicle categories.

\end{document}
